\documentclass[conference]{IEEEtran}
\IEEEoverridecommandlockouts
\usepackage{cite}
\usepackage{amsmath,amssymb,amsfonts}
\usepackage{graphicx}
\usepackage{textcomp}
\usepackage{xcolor}

\usepackage{subcaption}

\usepackage[draft,hidelinks,bookmarks=false,bookmarksdepth=0]{hyperref}
\usepackage{xurl}

\usepackage{listings}

\def\BibTeX{{\rm B\kern-.05em{\sc i\kern-.025em b}\kern-.08em
    T\kern-.1667em\lower.7ex\hbox{E}\kern-.125emX}}
\begin{document}

\title{Carbon Footprint of Relational and Hybrid Object-Relational Storage for Image-Based Time-Series Workloads}

\author{
\IEEEauthorblockN{Alfa Yohannis\textsuperscript{1}, Alexander Waworuntu\textsuperscript{2}, Master Edison Siregar\textsuperscript{1}, Arya Sanjaya\textsuperscript{1}, Refgiufi Avrianto\textsuperscript{1}}
\IEEEauthorblockA{\textsuperscript{1}Universitas Pradita, Indonesia\\
\{alfa.ryano, edison.siregar, arya.sanjaya, refgiufi.patria\}@pradita.ac.id}
\IEEEauthorblockA{\textsuperscript{2}Universitas Multimedia Nusantara, Indonesia\\
alex.wawo@umn.ac.id}
}


\maketitle

\begin{abstract}
This paper compares two architectures for image-based time-series workloads in smart building IoT: PostgreSQL with TimescaleDB storing images inline as \texttt{BYTEA} (PG), and PostgreSQL with TimescaleDB externalising binaries to MinIO object storage (PM). Using isolated Docker deployments and hardware-level energy tracking via CodeCarbon, we benchmark seven quality-90 JPEG resolutions from 360p to 5K. For insertion, PG is faster from 360p through 1440p (up to $\sim$$8\times$); PM becomes faster at 4K--5K. For point reads, PG leads only at 360p--480p; PM is faster from 720p onward ($\sim$$8\times$ at 5K). At the whole-phase level, PM reduces energy and CO$_2$eq emissions by $\sim$27\%. Profile-level estimates show PM emitting more at small resolutions but less from 1440p upward. These findings suggest that inline \texttt{BYTEA} suits lower-resolution ingest-heavy workloads, whereas externalised object storage becomes more suitable as payload size and carbon efficiency grow in importance.
\end{abstract}

\begin{IEEEkeywords}
time-series databases, IoT building monitoring, PostgreSQL, MinIO, object storage, carbon footprint, green computing
\end{IEEEkeywords}

\section{Introduction}
\label{sec:introduction}

Smart building and facility-monitoring deployments increasingly combine conventional telemetry with image streams from embedded cameras, doorbells, inspection devices, and other edge systems~\cite{ieee_ce_definition,intuz_iot_databases}. These applications often need to retain binary media together with timestamps, device identifiers, and other metadata so that operators can inspect the latest frame, replay recent events, or query activity over time~\cite{ceec_tsdb_review,jensen_tsms_survey}. In such systems, the data tier runs continuously and its energy use contributes directly to the operational carbon footprint of the building or edge infrastructure~\cite{nerc_tsdb_survey}.

Many organisations prefer to build on PostgreSQL~\cite{postgresql_official} and TimescaleDB~\cite{timescaledb_official} rather than introduce a specialised storage stack. When binary media is small, inline storage inside relational rows can be operationally convenient. When binary media becomes large, however, a common alternative is to keep only queryable metadata in PostgreSQL and store the bulk payload in an S3-compatible object store such as MinIO~\cite{minio_official} or Amazon S3~\cite{aws_s3_official}. These two options imply different write amplification, memory behaviour, and retrieval paths, and therefore potentially different energy profiles.

Existing benchmarks in the time-series literature focus primarily on lightweight numeric events~\cite{tsbs,cratedb_ts_benchmark,diva_timescaledb_mongodb_eval}. They rarely evaluate multi-megabyte image payloads, and even more rarely report energy or CO$_2$-equivalent emissions. This leaves a practical gap for image-heavy IoT deployments that must decide whether the simplicity of inline database storage outweighs the performance and sustainability benefits of separating metadata from binary objects.

This paper addresses the following research question: \textit{What are the performance and carbon-footprint trade-offs of inline PostgreSQL \texttt{BYTEA} storage (PG) versus PostgreSQL+MinIO externalised object storage (PM) for image-based time-series workloads across increasing image resolutions?}

To answer this question, we benchmark both architectures across seven JPEG profiles from 360p to 5K in isolated Docker containers, measure whole-phase energy and emissions with CodeCarbon~\cite{codecarbon}, and analyse insertion throughput, latest-payload point-read latency, storage amplification, and per-profile carbon estimates.

\section{Related Work}
\label{sec:related_work}

Time-series databases are designed to sustain append-oriented ingestion by combining time-based partitioning, compression, and storage layouts that reduce the write amplification of traditional B-Tree update paths~\cite{dewaal_timeseries_survey,jensen_tsms_survey}. TimescaleDB~\cite{timescaledb_official} extends PostgreSQL~\cite{postgresql_official} with hypertables while preserving SQL semantics and the operational ecosystem many organisations already use.

Most comparative benchmarks in this space target lightweight scalar data rather than binary media. TSBS~\cite{tsbs} uses small DevOps and telemetry records, CrateDB's benchmark~\cite{cratedb_ts_benchmark} focuses on numeric IoT metrics, and the evaluation by Mohamud and Brostr\"om~\cite{diva_timescaledb_mongodb_eval} studies numerical sensor workloads. These benchmarks do not characterise the behaviour of multi-megabyte image payloads, where storage-engine internals become far more consequential.

That gap matters because large binaries stress systems differently. PostgreSQL's TOAST~\cite{postgresql_toast_docs} mechanism relocates oversized values out of the main row, but the payload still remains within the database's transactional and retrieval machinery. Object stores such as MinIO~\cite{minio_official}, by contrast, are designed to stream large immutable objects directly while leaving relational systems responsible only for lightweight metadata and indexing. The architectural trade-off is therefore not simply ``one system versus two systems,'' but ``one transactional path for everything'' versus ``separating metadata management from binary transport.''

The carbon footprint of software systems has also attracted increasing attention. Lannelongue et al.~\cite{lannelongue_green_algorithms} formalised the Green Algorithms approach, and CodeCarbon~\cite{codecarbon} provides a practical implementation that estimates energy use and associated emissions from machine-level measurements and regional grid intensity. Within database systems, however, relatively little work applies this style of accounting to binary time-series storage workloads. Our study addresses that gap for relational versus hybrid object-relational image storage.

\section{Methodology}
\label{sec:methodology}

\subsection{Hardware and Software Environment}
All experiments execute on a single host machine running Ubuntu~24.04.4~LTS with Linux kernel 6.8.0-107-generic. The machine provides an 11th Gen Intel Core i9-11900H processor (8 cores / 16 threads, 2.50~GHz base, 4.90~GHz max), 30.99~GB RAM, and a 1.9~TB Samsung MZVL22T0HBLB-00BH1 NVMe SSD. The machine is located in Indonesia (Banten grid region); CodeCarbon uses the corresponding national carbon intensity factor. CPU and RAM power are measured via Intel RAPL (Running Average Power Limit) hardware counters (\texttt{tracking\_mode=machine}). Containers share host CPU and memory with no explicit CPU pinning or memory limits. The software stack consists of Docker~29.4.0 with Docker Compose~v5.1.1, PostgreSQL~15 with TimescaleDB (\texttt{latest-pg15}), MinIO (\texttt{minio/minio:latest}), Python~3.14.3, and the benchmark libraries \texttt{psycopg2}~2.9.11, \texttt{minio}~7.2.20, Pillow~12.2.0, and CodeCarbon~3.2.6.

\subsection{Environment and Workloads}
Both architectures are deployed in Docker containers on the same host. PostgreSQL exposes port~55432 and is launched with \texttt{max\_connections=100}. MinIO exposes ports~59000 and~59001 and uses the bucket \texttt{sensor-media}. The benchmark harness generates deterministic JPEG payloads from a single source image using a collage transform implemented in Python with Pillow. All profiles use JPEG quality~90.

Seven image profiles are evaluated: 360p ($640{\times}360$, $\sim$0.09~MB), 480p ($854{\times}480$, $\sim$0.16~MB), 720p ($1280{\times}720$, $\sim$0.36~MB), 1080p ($1920{\times}1080$, $\sim$0.80~MB), 1440p ($2560{\times}1440$, $\sim$1.40~MB), 4K~UHD ($3840{\times}2160$, $\sim$2.99~MB), and 5K ($5120{\times}2880$, $\sim$5.09~MB). For each profile, the benchmark performs a 100-row warm-up, followed by 5 measured insertion runs of 2,000 rows each in batches of 50. Latest-payload retrieval is measured as a point-read query with 2 warm-up runs and 5 measured runs. Driver overhead uses 10 warm-up and 10 measured roundtrips. The boxplot figures (Figures~\ref{fig:insert_throughput}--\ref{fig:retrieval_latency}) display the interquartile range and median across all runs; observed variance remained low, supporting the adequacy of 5 runs.

A representative PG row at 1080p contains: \texttt{sample\_id}=1, \texttt{device\_id}=1, \texttt{ts}=2025-04-20T00:00:00Z, \texttt{payload\_data}=$\sim$0.80~MB JPEG binary, \texttt{payload\_size\_bytes}=838860, \texttt{width}=1920, \texttt{height}=1080. For PM, the same row replaces \texttt{payload\_data} with \texttt{minio\_object\_key}=``sensor-media/1080p/frame\_0001.jpg''.

\subsection{Schema and Operations}
PG stores each frame as a TimescaleDB hypertable row with a \texttt{BYTEA} payload column, allowing PostgreSQL's TOAST~\cite{postgresql_toast_docs} mechanism to manage oversized values transparently. PM stores the metadata in PostgreSQL but replaces the inline payload with a \texttt{minio\_object\_key}; the binary image itself is uploaded to MinIO. Both tables use a compound index on \texttt{(device\_id, ts DESC)} to support latest-payload access.

Each measured insert run recreates the target table, and PM also recreates the MinIO bucket. For retrieval, PG returns the inline payload through PostgreSQL, whereas PM first retrieves the latest object key from PostgreSQL and then issues an HTTP GET to MinIO. This design isolates the core storage behaviour of each architecture while preserving realistic end-to-end operations.

\subsection{Carbon Footprint Measurement}
Each architecture phase runs as a single CodeCarbon~\cite{codecarbon}-tracked process. CodeCarbon reads CPU and RAM power via Intel RAPL hardware counters and converts total energy to CO$_2$-equivalent using the regional grid intensity for Indonesia ($\approx$0.676~kg CO$_2$eq/kWh). Each tracked phase therefore includes driver overhead, warm-up work, insert benchmarks, point-read benchmarks, and setup/teardown costs inside one measured duration. CodeCarbon reports total phase duration, energy consumed, CPU energy, RAM energy, and CO$_2$-equivalent emissions.

To obtain per-profile carbon estimates, we apply a constant power-rate model to the measured phase totals. For a phase with total energy $E$, emissions $C$, and duration $T$, we derive average rates $E/T$ and $C/T$ and multiply them by each profile's measured insert and point-read durations. These per-profile values are estimates rather than direct measurements, while the phase-level totals remain directly measured by CodeCarbon.

\begin{table}[htbp]
\centering
\renewcommand{\arraystretch}{1.1}
\caption{Insert Throughput, Point-Read Latency, and Storage Amplification (dataset size: 2,000 rows)}
\label{tab:performance_summary}
\begin{tabular}{l rr rr rr}
\hline
\textbf{Res.} & \multicolumn{2}{c}{\textbf{Insert (rows/s)}} & \multicolumn{2}{c}{\textbf{Point Read (ms)}} & \multicolumn{2}{c}{\textbf{Storage amp.}} \\
 & \textbf{PG} & \textbf{PM} & \textbf{PG} & \textbf{PM} & \textbf{PG} & \textbf{PM} \\
\hline
360p  & \textbf{738.32} & 104.05 & \textbf{0.720} & 1.717 & 1.0453 & \textbf{1.0029} \\
480p  & \textbf{475.75} & 56.88 & \textbf{1.017} & 1.163 & 1.0435 & \textbf{1.0017} \\
720p  & \textbf{205.12} & 51.69 & 2.589 & \textbf{1.767} & 1.0394 & \textbf{1.0007} \\
1080p & \textbf{91.86} & 39.02 & 3.641 & \textbf{2.352} & 1.0388 & \textbf{1.0003} \\
1440p & \textbf{50.51} & 35.47 & 7.909 & \textbf{2.280} & 1.0390 & \textbf{1.0002} \\
4K    & 23.04 & \textbf{28.49} & 14.942 & \textbf{2.461} & 1.0378 & \textbf{1.0001} \\
5K    & 13.16 & \textbf{15.59} & 27.322 & \textbf{3.360} & 1.0376 & \textbf{1.0001} \\
\hline
\end{tabular}

\footnotesize PG = PostgreSQL+TimescaleDB with inline \texttt{BYTEA}. PM = PostgreSQL+TimescaleDB+MinIO.
\end{table}

\begin{table*}[htbp]
\centering
\renewcommand{\arraystretch}{1.1}
\caption{Per-Profile Carbon Footprint Estimates from the Constant Power-Rate Model}
\label{tab:carbon_summary}
\begin{tabular}{l rr rr rr rr rr}
\hline
\textbf{Res.}
  & \multicolumn{2}{c}{\textbf{Duration (s)}}
  & \multicolumn{2}{c}{\textbf{Energy (\(\mu\)Wh)}}
  & \multicolumn{2}{c}{\textbf{CPU (\(\mu\)Wh)}}
  & \multicolumn{2}{c}{\textbf{RAM (\(\mu\)Wh)}}
  & \multicolumn{2}{c}{\textbf{Emissions (mg CO$_2$eq)}} \\
  & PG & PM & PG & PM & PG & PM & PG & PM & PG & PM \\
\hline
360p  & \textbf{17.3} & 103.6 & \textbf{218.4} & 979.3 & \textbf{122.5} & 403.9 & \textbf{96.0} & 575.4 & \textbf{147.6} & 662.0 \\
480p  & \textbf{27.4} & 183.9 & \textbf{346.7} & 1738.4 & \textbf{194.4} & 716.9 & \textbf{152.3} & 1021.5 & \textbf{234.3} & 1175.0 \\
720p  & \textbf{62.3} & 202.3 & \textbf{788.3} & 1912.6 & \textbf{442.0} & 788.8 & \textbf{346.3} & 1123.8 & \textbf{532.8} & 1292.8 \\
1080p & \textbf{141.6} & 273.8 & \textbf{1790.7} & 2588.8 & \textbf{1004.1} & 1067.7 & \textbf{786.6} & 1521.2 & \textbf{1210.4} & 1749.9 \\
1440p & \textbf{254.3} & 301.2 & 3216.1 & \textbf{2847.6} & 1803.3 & \textbf{1174.4} & \textbf{1412.7} & 1673.2 & 2173.8 & \textbf{1924.8} \\
4K    & 564.6 & \textbf{371.8} & 7140.1 & \textbf{3514.8} & 4003.7 & \textbf{1449.6} & 3136.5 & \textbf{2065.3} & 4826.2 & \textbf{2375.8} \\
5K    & 993.0 & \textbf{671.1} & 12558.4 & \textbf{6344.6} & 7041.8 & \textbf{2616.6} & 5516.6 & \textbf{3728.0} & 8488.6 & \textbf{4288.5} \\
\hline
\end{tabular}
\vspace{2pt}
\begin{minipage}{\textwidth}
\footnotesize Measured whole-phase totals including driver overhead, warm-up, and setup: PG = 0.029330~kWh / 19824.8~mg CO$_2$eq; PM = 0.021427~kWh / 14483.0~mg CO$_2$eq.
\end{minipage}
\end{table*}

\section{Results}
\label{sec:results}

\begin{figure}[t]
\centering
\includegraphics[width=\columnwidth]{figures/boxplot_insert_throughput.pdf}
\caption{Insert throughput across 360p--5K. PG leads through 1440p, while PM overtakes at 4K and remains ahead at 5K.}
\label{fig:insert_throughput}
\end{figure}

\begin{figure}[t]
\centering
\includegraphics[width=\columnwidth]{figures/boxplot_retrieval_latency.pdf}
\caption{Latest-payload point-read latency across 360p--5K. PG is slightly faster at 360p and 480p, while PM is faster from 720p onward and the gap widens sharply at 4K and 5K.}
\label{fig:retrieval_latency}
\end{figure}

\begin{figure*}[t]
\centering
\includegraphics[width=.9\textwidth]{figures/carbon_footprint.pdf}
\caption{Estimated per-profile CO$_2$eq emissions (mg) split into insert and point-read contributions. PM emits more at small profiles but less from 1440p upward.}
\label{fig:carbon_footprint}
\end{figure*}

\begin{figure*}[t]
\centering
\includegraphics[width=.9\textwidth]{figures/carbon_breakdown.pdf}
\caption{Per-profile energy breakdown (CPU + RAM, stacked bars) with benchmark duration shown as dashed lines. PM's carbon crossover is consistent with the crossover in end-to-end benchmark time.}
\label{fig:carbon_breakdown}
\end{figure*}

\subsection{Supporting Performance Observations}
\label{subsec:performance}
Table~\ref{tab:performance_summary} and Figures~\ref{fig:insert_throughput}--\ref{fig:retrieval_latency} show a write-side crossover. PG is faster from 360p through 1440p, with especially large margins at 360p (738.32 vs.\ 104.05~rows/s, $7.10\times$) and 480p (475.75 vs.\ 56.88~rows/s, $8.36\times$). The gap narrows steadily as payload size increases and reverses at 4K, where PM reaches 28.49 vs.\ 23.04~rows/s ($1.24\times$ faster), and remains in PM's favour at 5K (15.59 vs.\ 13.16~rows/s, $1.18\times$ faster). On this host, inline \texttt{BYTEA} provides higher insert throughput up to 1440p, while PM becomes faster at the two largest profiles.

Latest-payload point-read latency follows a different pattern. PG is slightly faster at the two smallest profiles: 0.720 vs.\ 1.717~ms at 360p and 1.017 vs.\ 1.163~ms at 480p. From 720p onward, PM becomes faster. At 1080p, PM reduces latency from 3.641 to 2.352~ms ($1.55\times$ faster); at 4K, from 14.942 to 2.461~ms ($6.07\times$ faster); and at 5K, from 27.322 to 3.360~ms ($8.13\times$ faster). The read-side crossover therefore occurs earlier than the write-side crossover, but only after the two smallest profiles.

PM consistently shows lower storage amplification. PG remains near $1.04\times$ storage amplification for all profiles, whereas PM stays essentially at raw-payload size, between $1.0029\times$ and $1.0001\times$. At 5K and 2,000 inserted rows, PG occupies 10555.0~MB versus 10172.9~MB for PM, a 382.1~MB gap despite similar logical content.

Driver overhead provides additional context. Averaged across the seven profiles, PostgreSQL's driver roundtrip measures 0.062~ms, while the MinIO \texttt{list\_buckets} roundtrip used here measures 0.606~ms. This auxiliary measurement suggests a higher fixed client-side roundtrip on the MinIO path, although it does not by itself determine the end-to-end insert or point-read results.

\subsection{Carbon Footprint (RQ)}
\label{subsec:carbon_footprint}
The whole-phase CodeCarbon measurements show that PM is more carbon-efficient overall on this seven-profile benchmark. The PostgreSQL-only phase ran for 2319.1~s, consumed 0.029330~kWh, and emitted 19824.8~mg CO$_2$eq. The PostgreSQL+MinIO phase ran for 2266.3~s, consumed 0.021427~kWh, and emitted 14483.0~mg CO$_2$eq. This corresponds to a \textbf{26.9\% reduction} in both energy and emissions for PM at the phase level.

Table~\ref{tab:carbon_summary} and Figures~\ref{fig:carbon_footprint}--\ref{fig:carbon_breakdown} refine that picture by estimating the profile-level contribution of insert and point-read work. These estimates show that PM is \emph{not} uniformly greener. At 360p, PM emits 662.0~mg CO$_2$eq versus 147.6~mg for PG (+348.3\%). At 480p, the PM penalty is even larger (1175.0 vs.\ 234.3~mg, +401.4\%), reflecting its poor small-payload insert throughput. PM remains worse at 720p and 1080p, though the gap narrows substantially by 1080p (1749.9 vs.\ 1210.4~mg, +44.6\%).

The carbon crossover occurs at 1440p. There, PM emits 1924.8~mg versus 2173.8~mg for PG, an 11.5\% advantage. At 4K the PM advantage becomes pronounced, dropping from 4826.2 to 2375.8~mg (-50.8\%), and at 5K the reduction remains substantial at 4288.5 versus 8488.6~mg (-49.5\%). This pattern is consistent with the throughput crossover observed in Section~\ref{subsec:performance}: once payload sizes become large enough for PostgreSQL's inline binary handling to dominate execution time, PM's separate object-storage path is associated with both faster completion and lower estimated carbon.

The difference between the measured whole-phase totals and the summed profile estimates is expected. The per-profile table only allocates insert and point-read time, whereas the tracked phases also include warm-up activity, driver-overhead measurements, and setup work. The directly measured phase totals therefore remain the correct basis for the overall RQ answer, while the profile estimates explain where that overall difference originates.

\section{Discussion}
\label{sec:discussion}

\textbf{Resolution-dependent crossover.} The updated benchmark does not support a universal ``always use MinIO'' or ``always use inline \texttt{BYTEA}'' conclusion. PG remains faster for writes through 1440p on this machine and also keeps a small point-read advantage at 360p and 480p. From 720p onward, however, PM becomes the faster point-read path; from 4K onward it also becomes the faster write path; and from 1440p onward it becomes the lower-carbon option.

\textbf{Why PostgreSQL leads at smaller writes.} At lower resolutions, PG benefits from batched inserts that avoid an extra object-store request per payload. The images are small enough that PostgreSQL's inline binary handling and TOAST overhead remain tolerable, while PM still pays the fixed cost of a MinIO PUT plus a relational metadata write. That is why PM performs particularly poorly at 360p and 480p despite its much better storage amplification.

\textbf{Why PostgreSQL+MinIO becomes faster at large payloads.} As payload size grows, the cost of keeping the binary within the relational engine rises sharply. PostgreSQL must still log, manage, and retrieve large values through the database path, whereas PM lets PostgreSQL manage only metadata while MinIO streams the object directly. This architectural separation is consistent with the observed performance and carbon advantages once the binary payload is large enough, which in our measurements happens around 1440p for emissions and 4K for insertion throughput.

\textbf{Implications for system design.} These results suggest a workload-sensitive deployment heuristic. If the system is dominated by lower-resolution ingestion and operational simplicity is a priority, inline \texttt{BYTEA} storage may remain practical. If the system needs lower-latency latest-frame access from 720p upward, or expects sustained 4K--5K payloads, externalising binaries to object storage may become the more suitable choice. PM also offers consistently better storage efficiency, which compounds as retention increases.

\textbf{Database selection rationale.} This study uses PostgreSQL as the sole relational baseline because TimescaleDB---the most widely adopted open-source time-series extension---is built exclusively on PostgreSQL. Other relational engines lack a comparable native time-series extension. The inline-vs-externalized trade-off is conceptually transferable to any relational database supporting large-object columns, though exact crossover points would depend on each engine's page layout and WAL implementation.

\textbf{Scalability considerations.} Our benchmark uses 2,000 rows per resolution. At larger scales, PG's inline BYTEA disadvantage at high resolutions should compound: TOAST tables grow, and WAL volume scales linearly with payload size. PM writes only lightweight metadata rows to PostgreSQL, so WAL and index overhead scale with metadata rather than payload size. MinIO's horizontal scale-out with erasure coding could further favour PM at large data volumes.

\textbf{Threats to validity.} Results are from a single host and a single container topology. Different hardware, distributed object-storage deployments, or more compressible media could shift the exact crossover points. The fixed batch size of 50 records may not reflect all production ingestion patterns; alternative batch sizes could shift the observed carbon crossover. CPU and RAM power are measured via Intel RAPL hardware counters (\texttt{tracking\_mode=machine}); the per-profile carbon table is a proportional estimate based on a constant power-rate assumption rather than a hardware-instrumented per-operation trace. All benchmarks use sequential single-client access; concurrent multi-client workloads, mixed read/write patterns, and diverse query types were not evaluated. Even so, both architectures were benchmarked under identical hardware conditions, making the comparative trends informative within this benchmark.

\section{Conclusions}
\label{sec:conclusions}

This study evaluated the performance and carbon footprint of two storage architectures for image-based time-series workloads: PostgreSQL with inline \texttt{BYTEA} storage and PostgreSQL with externalised MinIO object storage, using CodeCarbon with RAPL-based hardware energy measurement.

\textbf{RQ answer:} The better architecture depends on payload size and workload emphasis, but the complete seven-profile benchmark on this machine shows lower measured whole-phase energy use and emissions for PostgreSQL+MinIO overall. Measured whole-phase totals show PM emitting 14483.0~mg CO$_2$eq versus PostgreSQL's 19824.8~mg, a \textbf{26.9\% reduction}. PG remains faster for insertion from 360p through 1440p, while PM becomes faster at 4K and 5K. For latest-payload point reads, PG is slightly faster at 360p and 480p, but PM is faster from 720p onward and the gap reaches $8.1\times$ at 5K.

The key practical finding is therefore a set of crossover points rather than a universal winner. On this host, inline PostgreSQL may suit smaller ingest-heavy workloads, while PostgreSQL+MinIO may become more suitable as binary payload size increases and point-read or carbon efficiency become more important. Future work should evaluate the same architectures on multi-node deployments and under higher client concurrency to determine how stable these crossover points remain outside a single-host setup. Additionally, PostgreSQL Large Objects (LO) merit investigation as a third architecture---using \texttt{lo\_import}/\texttt{lo\_export} instead of inline BYTEA---to determine whether LO's streaming interface narrows the carbon and retrieval gap with MinIO.

\section*{Acknowledgment}
The authors thank Universitas Pradita and Universitas Multimedia Nusantara for supporting this research.


\bibliographystyle{IEEEtran}
\bibliography{references}

\end{document}
